% !TeX root = ../教材题库_全部条目_转写工作稿.tex
% Source: XJ-B1-C04-RV-P152-Y0954
% 第4章 / 章末复习 复习题四 / 教材 p.152 / PDF 第159页
% OCR未校对：公式与文字均可能有误；未完成数学TeX转写；禁止直接发给学生
% 原题及相邻原解截图（核对用）：../题目原图/XJ-B1-C04-RV-P152-Y0954-01.webp
% 本文件只在转写完成后作为题干使用；OCR 条目尚不是可靠题干。
\begin{bankitem}{复习题四 第17题}{第4章 / 章末复习 复习题四 / 教材 p.152 / PDF 第159页}{XJ-B1-C04-RV-P152-Y0954}{OCR未校对：公式与文字均可能有误；未完成数学TeX转写；禁止直接发给学生}
\begin{ocrdraft}
[教材p.152 / PDF第159页]
\par
17．煤油在作为喷气发动机燃料之前需通过呈管道状的黏土以去除其中的污染
\par
物．某种煤油中污染物含量为，测得这种煤油通过长为工m的管道后污染物的含
\par
量如下表：
\par
x/m
\par
0
\par
1
\par
2
\par
3
\par
炻
\par
0．8力0
\par
0．64
\par
0．512力0
\par
（1）写出一个关于的函数解析式，使之与表中数据吻合；
\par
（2）要使这种煤油中污染物的含量不超过原来的5％，利用（D中得到的函数求
\par
所要通过的管道至少要多长()g0．8
\par
一0，097，lg0．05一1.30D．
\end{ocrdraft}
\end{bankitem}
